On considère la fonction $f$ définie sur l'intervalle $]0~;~+\infty[$ par 

\[f(x) = x(\ln x)^2.\]

On admet que la fonction $f$ est dérivable sur l'intervalle $]0~;~+\infty[$. On note $f'$ sa fonction dérivée.

\medskip

\begin{enumerate}
\item Déterminer la limite de la fonction $f$ en $+\infty$.
\item Pour tout réel $x > 0$, on pose $g(x) = x\ln x$.
	\begin{enumerate}
		\item Démontrer que pour tout réel $x > 0$, on a $f(x) = 4\left(g\left(\sqrt x\right)\right)^2$.
		\item En déduire $\displaystyle\lim_{x \to 0} f(x)$.
	\end{enumerate}
\item Dans cette question, on étudie les variations de la fonction $f$ sur l'intervalle $]0~;~+\infty[$.
	\begin{enumerate}
		\item Démontrer que sur l'intervalle $]0~;~+\infty[$, \: $f'(x) = (\ln x)(2 + \ln x)$.
		\item En déduire les variations de la fonction $f$ sur l'intervalle $]0~;~+\infty[$.
		\item Donner la valeur exacte du maximum de la fonction $f$ sur l'intervalle $]0~;~1]$.
	\end{enumerate}
\item On considère l'équation $f (x) = 2$.
	\begin{enumerate}
		\item Justifier que, sur l'intervalle $]0~;~+\infty[$, cette équation admet une unique solution.

On note $\alpha$ cette solution.
		\item Donner un encadrement de $\alpha$ d'amplitude 0,1.
	\end{enumerate}
\item Soit $a$ un nombre réel appartenant à l'intervalle ]0~;~1].
	\begin{enumerate}
		\item Donner une interprétation géométrique de $\displaystyle\int_a^1 f(x)\: \text{d}x$.
		\item À l'aide d'une intégration par parties, justifier que : 

\[\displaystyle\int_a^1 f(x)\: \text{d}x = - \dfrac{a^2}{2}(\ln a)^2 - \displaystyle\int_a^1 x \ln x\: \text{d}x.\]

		\item En utilisant à nouveau une intégration par parties, démontrer que :

\[\displaystyle\int_a^1 f(x)\: \text{d}x =  - \dfrac{a^2}{2}(\ln a)^2 + \dfrac{a^2}{2} \ln a + \dfrac14 - \dfrac{a^2}{4}.\]

		\item Déterminer la limite de $\displaystyle\int_a^1 f(x)\: \text{d}x$ quand $a$ tend vers 0.
	\end{enumerate}
\end{enumerate}
